\scanpage{343}
% @ top centre
\mspage{67}

% @ top paragraph, lines 1–4
\deleted{déjà} $W$-\unsure{cocartésiens|cocartésien}. \deleted{\unsure{Pourtant}} Mais il me semble à présent que nous pouvons le prouver, \inserted{\unsure{i.e.}} sous la seule condition supplémentaire \deleted{\unsure{$W$}} $\mathrm{Loc}(3\mathrm{bis})^{*}$,
% @ upper left opposite asterisk, three diagonal lines
% slanted
\margin{et peut-être, de plus, $\mathrm{Loc}(1\ \mathrm{ter})$ $W$ fortement saturé}
% @ upper third, continuation to diagram pair
laquelle assure l'existence \unsure{des foncteurs|du foncteur} $f_!$, et en particulier la construction du \deleted{\unsure{carré}} $\mathrm{Holim}$ d'un diagramme

% @ upper third, two forks with interposed comme
\diagram{\begin{tikzcd}[ampersand replacement=\&,row sep=0pt,column sep=small]
 \& \& X_1 \& \& \& X_1\\
 \& X_0\arrow[ur]\arrow[dr]\& {} \& \text{comme}\quad\displaystyle\int \& X_0\arrow[ur]\arrow[dr]\& {}\\
 \deleted{\unsure{\int|f|\text{Or}}},\& \& X_{\unsure{2|\sim}}\& \& \& X_2.
\end{tikzcd}}
% @ centre, statement before open square
Donc dire qu'un diagramme de MV ouvert

% @ centre, open square in diamond position
\diagram{\begin{tikzcd}[ampersand replacement=\&,row sep=0pt,column sep=small]
 \& X_1\arrow[dr]\& {}\\
 X_0\arrow[ur]\arrow[dr]\& \& X\\
 \& X_2\arrow[ur,dash]\& {}
\end{tikzcd}}
% @ centre to lower, canonical map
est $W$-cocartésien, doit revenir à dire que la flèche canonique

% @ lower middle, integral of fork and map
\diagram{\begin{tikzcd}[ampersand replacement=\&,row sep=0pt,column sep=small]
 \& \& X_1\& {}\\
 \displaystyle\int\& X_0\arrow[ur]\arrow[dr]\& {}\arrow[r]\& X\\
 \& \& X_2\& {}
\end{tikzcd}}
% @ lower third, qualification and references
est dans $W$; \deleted{\unsure{du}} \unsure{moins|disons} si $W$ est supposé fortement saturé. Or cela a été bel et bien vu dans $\underline{\overline{\mathrm{XII}}}$ 3.4. Th.4 et son \uline{corollaire} 3, p.149). Dans loc cit. j'énonce les conclusions sous l'hypothèse que \inserted{\unsure{j'appelais}} $W(\unsure{A|7}\mathrm{bis})$ (l'actuel $\mathrm{Loc}(6)$), mais je viens de vérifier que pour la partie \uline{directe}, donc aussi pour ses cor. 1 à 3, on n'a pas besoin
% archive sheet number: 342
